\subsection{Nuova scrittura per gli stati}
\begin{align*}
\ket{\psi} &= \braket{+|\psi} \ket{+} + \braket{-|\psi} \ket{-} = \ket{+} \braket{+|\psi} \ket{-} \braket{-|\psi} = \\
&= \underbrace{\left( \ket{+} \bra{+} + \ket{-} \bra{-} \right)}_{\mathbbm{1}} \ket{\psi}
\end{align*}
In generale, gli stati \(\ket{e_i}\) tali che \(\braket{e_i|e_j} = \delta_{ij}\) danno
\[
\mathbbm{1} = \sum_i \ket{e_i}\bra{e_i} \qquad \text{relazione di completezza}
\]
Ad esempio,
\[
\braket{\varphi|\psi} = \sum_i \braket{\varphi|e_i} \braket{e_i|\psi}
\]
e allora
\[
\ket{\psi} = \sum_i \ket{e_i} \braket{e_i|\psi} = \sum_i \ket{e_i} c_i^{(\psi)}
\]
\[
\ket{\varphi} = \sum_i \ket{e_i} \braket{e_i|\varphi} = \sum_i \ket{e_i} c_i^{(\varphi)}
\]
\[
\braket{\varphi|\psi} = \sum_i \left( c_i^{(\varphi)} \right)^* c_i^{(\psi)}
\]
e quindi
\[
\braket{\psi|\psi} = \sum_i \braket{\psi|e_i}\braket{e_i|\psi} = \sum_i \left| c_i^{(\psi)} \right|^2
\]